\documentclass[./eval.tex]{subfiles}
\newcommand{\head}{\noindent Máquinas Eléctricas II |Evaluación| Apellidos y Nombres.
\rule{10cm}{0.1mm}}
\newcommand{\separador}{\noindent\rule{\paperwidth}{0.1mm}\newline}
\begin{document}
\head

 \begin{enumerate}
 \item Un conductor se desplaza a una velocidad lineal de $v=20m/s$ dentro de un campo magnético fijo de $B=5T$ de inducción. Determinar el valor de la f.e.m. inducida en el mismo si posee una longitud de $l=200cm$.
    \item A que velocidad se deberá mover un conductor de longitud $l=0.5m$ para que dentro de una inducción de campo magnético de $B=2T$ produzca una f.e.m de $e=5V$.
    \item Determinar la inducción de campo magnético $B$, que será necesaria para que un conductor de $l=1m$ de longitud que se mueve a $v=20m/s$ cortando las líneas del campo magnético produzca una f.e.m de $e=12V$.
    \item ¿Cuál la f.e.m de auto-inducción si el flujo magnético $\Phi$ que existe en el campo magnético producido por una bobina de una electroválvula si tiene esta un núcleo cilíndrico de diámetro $d=4cm$ y la inducción magnética en la misma crece de $B=0T$ a $B=1.5T$ en $\Delta t=0.3s$?
    \item Calcular el valor de la f.e.m de auto-inducción que se desarrollará una boina con coeficiente de auto-inducción de 50 milihernios si se aplica una corriente que crece desde cero hasta $I=30A$ en un tiempo de $\Delta t=0,1 seg.$
    \separador
    \end{enumerate}
 \head
    \begin{enumerate}
        
    \item Cuanto tiempo habrá que esperar para que la f.e.m de auto-inducción se reduzca a $e_{auto}=1V$ si se hace circular $I=10A$ por una bobina con un coeficiente de auto-inducción de $L=60 mH$  
    \item Cual es la f.e.m de auto inducción del problema anterior cuando pasa 1 milisegundo de tiempo $t=0,001seg$ y ¿Cuál es el valor de la fem de auto-inducción después de 1 minuto?
    \item A que velocidad se deberá mover un conductor de longitud $l=0.5m$ para que dentro de una inducción de campo magnético de $B=1T$ produzca una f.e.m de $e=5V$.
    \item Determinar la inducción de campo magnético $B$, que será necesaria para que un conductor de $l=1m$ de longitud que se mueve a $v=10m/s$ cortando las líneas del campo magnético produzca una f.e.m de $e=12V$.
    \item ¿Cuál la f.e.m de auto-inducción si el flujo magnético $\Phi$ que existe en el campo magnético producido por una bobina de una electroválvula si tiene esta un núcleo cilíndrico de diámetro $d=4cm$ y la inducción magnética en la misma crece de $B=0T$ a $B=1.6T$ en $\Delta t=0.2s$?
    
       \separador
 \end{enumerate}
\end{document} 
